\documentclass[12pt]{article}
\usepackage{lineno}
%\usepackage{ametsoc}
\usepackage{natbib}
\usepackage{a4wide}
\usepackage{longtable}
\usepackage{amsfonts}
\usepackage{amssymb}

\usepackage{color}
\usepackage[pdftex]{graphicx}
\usepackage[pdftex,colorlinks, citecolor=black,linkcolor=black,urlcolor=black]{hyperref}
\definecolor{darkgreen}{rgb}{0.1,0.4,0.1}
\definecolor{darkblue}{rgb}{0.1,0.1,0.3}



\newif\ifdetail
\global\detailfalse
%\global\detailtrue

\def\bem#1{\colorbox{green}{\em #1}}

\newcommand{\beq}  { \begin{eqnarray} }
\newcommand{\eeq}  { \end{eqnarray}}
\newcommand{\beeq}  { \begin{eqnarray*} }
\newcommand{\eeeq}  { \end{eqnarray*}}
\newcommand{\eq }  [1]{Eq.~(\ref{#1})}  % Eq. (#1)
\newcommand{\eqs }  [1]{Eqs.~(\ref{#1})}
\newcommand{\sect}  [1]{Sec.~(\ref{#1})}  % Eq. (#1)
\newcommand{\fig}  [1]{Fig.~\ref{#1}}   % Fig. #1
\renewcommand{\v}[1]{{\mbox{\boldmath$ #1 $}}}  % Vektor durch Fettdruck
\newcommand{\m}    [1]{{\bf  #1 }}              % Matrix in sans serif
\newcommand{\gz}[1] {{{d #1} \over {dz}}}
\newcommand{\pt}[1]{{\partial_t #1}}
\newcommand{\ptt}[1]{{\partial_{tt} #1}}
\newcommand{\pttt}[1]{{\partial_{ttt} #1}}
\newcommand{\pttx}[1]{{{\partial^3 #1} \over {\partial t^2 \partial x}}}
\newcommand{\ptty}[1]{{{\partial^3 #1} \over {\partial t^2 \partial y}}}
\newcommand{\ptx}[1]{{{\partial^2 #1} \over {\partial t \partial x}}}
\newcommand{\pty}[1]{{{\partial^2 #1} \over {\partial t \partial y}}}
\newcommand{\pz}[1]{{\partial_z #1}}
\newcommand{\pzz}[1]{{\partial_{zz} #1}}
\newcommand{\px}[1]{{\partial_x #1}}
\newcommand{\py}[1]{{\partial_y #1}}
\newcommand{\pxi}[1]{{\partial_i #1}}
\newcommand{\pxj}[1]{{\partial_j #1}}
\newcommand{\pxk}[1]{{\partial_k #1}}
\newcommand{\pb}[1]{{{\partial #1} \over {\partial b}}}
%\newcommand{\pxy}[1]{{{\partial^2 #1} \over {\partial x \partial y}}}
\newcommand{\pxy}[1]{\partial_{xy} #1}
\newcommand{\pyy}[1]{\partial_{yy} #1}
\newcommand{\pxx}[1]{\partial_{xx} #1}
\newcommand{\kx}{{\bf k} \times}
\newcommand{\dt}[1]{{{d #1} \over {d t}}}
\newcommand{\Dt}[1]{{{D #1} \over {D t}}}
\newcommand{\Div}[1]{ \nabla \cdot #1 }
\renewcommand{\div}[1]{ \nabla_{yz} \cdot #1 }
\newcommand{\Bar}[1]{ \overline{ #1} }
\newcommand{\rvec} [1]{
    \raisebox{-1.5ex}{$\stackrel{\textstyle #1}{\neg}$} }% rotierter Vektor
\newcommand{\nablq}   {\rvec\nabla}                    % rotated nabla operator
\newcommand{\Vec}[1]{\left(\begin{array}{c} #1 \end{array}\right) }
%\renewcommand{\mho}{\mathcal{W}}
\newcommand{\pn}[1]{{{\partial #1} \over {\partial n}}}
\newcommand{\pnn}[1]{{{\partial #1} \over {\partial m}}}
\newcommand{\ps}[1]{{{\partial #1} \over {\partial s}}}
\newcommand{\p}{{\partial}}
\newcommand{\vn}{{\v \nabla}}
\newcommand{\lk}{\left(}
\newcommand{\rk}{\right)}

\newcommand{\tl}{\tilde}

\newcommand{\I}{\mathcal{I}}
\newcommand{\E}{\mathcal{E}}
\newcommand{\A}{\mathcal{A}}

\newcommand{\ep}{\epsilon}
\newcommand{\erfc}{\mathrm{erfc}}
 
 
\begin{document}

 \section*{The spectral wave model wim\_simple (Wave Interaction Model, WIM)}
 
wim\_simple is a spectral surface wave model, coupled to an ocean model in an energetically
and dynamically consistent way.
Basis is the wave action spectral density equation for deep water surface gravity waves
\beq \label{eqA}
 \pt{A}  +  \dot{\v x} \cdot  \vn_{\v x}  A   +\dot{\v k} \cdot  \vn_{\v k} A  &=& S = S_{in} + S_{out} + S_{nl}
\eeq
with wave action density $A(\v k,\v x,t)$, the horizontal wavenumber vector $\v k = (k_x,k_y)$,
the dispersion relation $\Omega=\sigma + \v U \cdot \v k$ with intrinsic frequency
$\sigma(\v k) = \sqrt{ g |\v k| }$ and Doppler shift velocity $\v U$,
transport velocities $  \dot{\v x} =   \vn_{\v k} \Omega$ and $\dot{\v k}= - \vn_{\v x} \Omega$,
action input by wave growth $S_{in}$, action dissipation by wave breaking $S_{out}$,
 and non-linear action transfer $S_{nl}$
and operators $\vn_{\v x} = (\p_x,\p_y)$ and $\vn_{\v k} = (\p_{k_x}, \p_{k_y})$.
Only deep water waves are considered, with $\vn_{\v x} \sigma = 0$ and $\p_t \sigma = 0$,
but the concept extends readily to waves at finite depth.
Since the scattering integral $S_{nl}$ conserves energy and momentum of the waves,
it holds that $\int \sigma S_{nl} d \v k = 0$ and $\int \v k S_{nl} d \v k = 0$.


The surface wave model predicts the Stokes velocity $\v u_S$ and the wave stress $\v \tau_o$
to force an ocean model.
Considering the Lagrangian ocean velocity $\v u_L = \v u + \v u_S$ and its kinetic energy,
 the consistent energy cycle becomes 
\beq \label{eq_K}
\p_t \int_{-\infty}^0  |\v u_L|^2/2 dz + \vn_{\v x}  \cdot \int_{-\infty}^0  \v u^L_h  \pi  dz  &=&   C  \\  \label{eq_W}
\p_t  \int E d \v k +  \vn_{\v x}  \cdot  \int \lk  \dot{\v x}   E  + \dot{ \v x}  A (\v U \cdot \v k)  \rk d \v k
   &=& \int (\Omega S_{in} +\v U \cdot \v k \, S_{nl}) d \v k +\int \sigma S_{out}   d \v k   -C  \nonumber
\eeq
with the wave energy spectral density $E = \sigma A$ and the exchange term 
\beq \label{eq_C}
   C&=&  \int_{-\infty}^0  \v u^L_h \cdot   \p_t \v u_S  dz +\int_{-\infty}^0  \v u^L_h  \cdot    \p_z \v \tau  dz 
   = \int  \v U \cdot \v k (\p_t A -S_{out}) d \v k
   \eeq
   and the setting, following Vanneste and Young for the Doppler shift velocity $\v U$,
\beq \label{eq_Q}
  \v U(\v x,\v k,t) = \int_{-\infty}^0 \v u^L_h Q  dz  ~~,~~\v u^S_h(\v x,z,t)  = \int Q A \v k d \v k~~,~~Q(k,z) = 2 k e^{2 kz}    
  \\
\p_z \v \tau(\v x,z,t) = -  \int  Q    \v k S_{out} d \v k~,~ \v \tau_o(\v x,t) =\int_{-\infty}^0 \p_z \v \tau dz =- \int   \v k S_{out} d \v k
\eeq
with $\int_{-\infty}^0 Q \, dz = 1$ and the ocean momentum forcing $ \v \tau_o =-\int   \v k S_{out} d \v k$,
i.e.\ the stress on the ocean is the dissipation of wave momentum.
Note that $\v U$ contains both the Eulerian mean and the Stokes horizontal velocity components.
Waves and ocean exchange energy by the effect of the Doppler shift and the ocean stress, both
given by $C$; other source terms of the wave energy are the wave growth by the wind,
$\int \Omega S_{in} d \v k$, and the dissipation by wave breaking, $\int \sigma S_{out} d \v k$,
which may be directed into a turbulent ocean energy reservoir (see below).
The cycle closes exactly up to the term $\int \v U \cdot \v k \, S_{nl} d \v k$, which represents
the work done on the sheared mean flow by the vertical redistribution of Stokes momentum
through the nonlinear transfer. Only the baroclinic part of $\v U$ contributes to this term,
since $\int \v k S_{nl} d \v k = 0$. The discrete interaction approximation used for $S_{nl}$
evaluates resonant quadruplets based on the intrinsic frequency and thus enforces
$\int \sigma S_{nl} d \v k = 0$, such that with a background current this term appears as a
spurious source; since it is small compared to the other sources, it is ignored.


\subsection*{Relation to WAM}

The concept and discretisation of wim\_simple follows closely the example of WAM
found in WAM (cycle 7) at \url{https://github.com/myWAveModel/WAM},
except for simplified action sources.
Instead of the canonical wavenumber coordinates $\v k$ used above,
frequency $\nu = \sigma/(2 \pi)$ and wavenumber angle $\theta$
with $k_x = k \sin \theta$, $k_y = k \cos \theta$ and $k=|\v k|$ are used as phase dimensions,
covered by an equidistant grid in $\theta$ and a logarithmic grid in $\nu$ as in WAM.
Different to WAM, however, $\nu$ is here the \emph{intrinsic} frequency, while in WAM
the extrinsic frequency is used; the intrinsic frequency is better suited to
represent wave--mean flow interaction, following the example of WAVEWATCH III.
Instead of the spectral energy density $E(\v k) = \sigma A(\v k)$, the model carries
the WAM-like spectral density $E(\theta,\nu)$, related by
$ \int_{-\infty}^0  |\v u_L|^2/2 dz= \int E(\v k) d \v k = g \int E(\theta,\nu) d \theta d\nu  $.
Units of $E(\v k) d \v k$ are $m^3/s^2$,
thus $E(\theta,\nu)  d \theta d\nu$ has units $m^2$ (the surface elevation variance)
and $E(\theta,\nu)$ has units $m^2 s $,
while $E(\v k)$ has units $m^5/s^2$; the factor $g$ shows up accordingly in the
definitions of $\v u_S$ and $\v \tau$ below.
\ifdetail
Coordinate transformation is not necessary but given here for completeness.
 From $(k_x,k_y)$ to ($\theta,k$) 
with $k_x=k \cos \theta$, $k_y=k \sin \theta$
\beq
E(\v k) d \v k  = k E(\v k) d \theta d k
\eeq
since (take absolute value)
\beq
\left |   \begin{array}{cc}
\p k_x /\p \theta & \p k_x/ \p k  \\
\p k_y /\p \theta & \p k_y/ \p k  \\
\end{array} \right|  = 
\left |   \begin{array}{cc}
-k \sin \theta & \cos \theta  \\
k \cos \theta &  \sin \theta\\
\end{array} \right| 
= - k \sin^2 \theta - k \cos^2 \theta = -k
\eeq
 Transformation from  $(\theta,k)$ to  $(\theta,\sigma)$ with $k = \sigma^2/g$ by
\beq
   d \theta d k =  d \theta \frac{\p k}{\p \sigma} d \sigma = 2 (\sigma/g)  d \theta  d \sigma
   ~~,~~
   d \v k =   2 (\sigma^3/g^2)  d \theta  d \sigma
\eeq 
\fi




The model equation formulated for $A(\theta,\nu)$ reads
\beq \label{eq6}
% \pt{ A}    +  \dot{\v x} \cdot  \vn_{\v x}  A  +\dot{\theta} \p_{\theta} A   +\dot{\sigma} \p_{\sigma} A
% &=& S_{in} + S_{out} + S_{nl} 
% \\
  \pt{ A}    +  \vn_{\v x}  \cdot   (\dot{\v x} A ) +\p_{\theta}  (\dot{\theta} A  )
 + \p_{\nu} ( \dot{\nu} A)
 &=&S_{in} + S_{out} + S_{nl} 
\eeq
We need expressions for group velocity $\dot{\v x}$ and the drift velocities $\dot{\theta}$
and $\dot \nu$.

\subsubsection*{Group velocity}

Since $\v U = \v U(k)$, the group velocity is given by
\ifdetail
\beq
 \dot x &=& \p_{k_x} \Omega = \p_{k_x} \sigma  + \p_{k_x}  Uk_x  + \p_{k_x}  V k_y  
 =   \sqrt{g} \p_{k_x}  (k_x^2 + k_y^2)^{1/4}+ U + k_x \p_{k_x}  U + k_y\p_{k_x}  V 
 \\ &=&  (c /2)  k_x/k + U + \v k \cdot \p_{k_x}  \v U 
 \\
 \dot y &=& \p_{k_y} \Omega = (c /2)  k_y/k+ \p_{k_y}  Uk_x  + \p_{k_y}  V k_y  
 =  (c /2)  k_y/k+ V + \v k \cdot \p_{k_y}  \v U   
 \eeq
 \fi
 \beq
 \dot{\v x} =  (c /2)  \v s + \v U + \v k \cdot \vn_{\v k}  \v U 
 \eeq
with the phase velocity  $c = \sigma/k = \sqrt{g} /\sqrt{k} = g /\sigma $
and the wave direction vector $\v s = (\sin \theta,\cos \theta)$.
\ifdetail
Note that
\beq 
\vn_{\v k}  ( \v  k \cdot    \v U)   = \Vec{  \p_{k_x} (U k_x + V k_y) \\ \p_{k_y} (U k_x + V k_y) }
= \v U + \Vec{   k_x  \p_{k_x} U  + k_y   \p_{k_x}  V \\  k_x \p_{k_y} U  +k_y  \p_{k_y} V   }
\\
= 
 \v U + \Vec{   (k_x  \p_{k} U  + k_y    \p_{k}  V) k_x/k  \\  
 (k_x \p_{k} U  +k_y  \p_{k} V) k_y/k    }
\eeq
with $ \p_{k_x} k  = % \p_{k_x} (k_x^2+k_y^2)^{1/2} = 1/2  (k_x^2+k_y^2)^{-1/2}  2 k_x = 
k_x/k$
and  $\v U = \v U(k)$
\fi
Since the wavenumber gradient of the Doppler velocity $\v U$
is purely radial (it points along $\v k$),
\beq \label{eq:kgradkU}
\v k \cdot \vn_{\v k} \v U
= \Vec{ (k_x \p_k U + k_y \p_k V)\, k_x/k \\
        (k_x \p_k U + k_y \p_k V)\, k_y/k }
= \lk \v s \cdot \p_k \v U \rk \, \v k
= k \lk \sin\theta\, \p_k U + \cos\theta\, \p_k V \rk \, \v s ,
\eeq
with  $\p_k \v U = (\p_k U, \p_k V)$. The full group velocity reads
\beq \label{eq:xdot_full}
\dot{\v x} = (c/2)\,\v s + \v U + k \lk \sin\theta\, \p_k U + \cos\theta\, \p_k V \rk \, \v s .
\eeq
With $k = (2\pi \nu)^2/g$ one has
$k\,\p_k U_i = k \p_\nu U_i \p \nu/\p k = \frac{\nu}{2}\,\p_\nu U_i$
%with $2 \pi \nu =  \sqrt{g k} $ and $  (2 \pi) \p \nu/\p k = g/2 /\sqrt{g k}  $ 
%and $k   (2 \pi) \p \nu/\p k = \sqrt{g k} /2 =\pi f $
%
 so the radial correction can be written with
a single scalar factor
\beq \label{eq:xdot_disc}
R = \frac{\nu}{2}\lk \v s \cdot \p_\nu \v U  \rk ,
\qquad
\dot{\v x} %= \Vec{ (c/2)\sin\theta + U + R\,\sin\theta  \\ (c/2)\cos\theta + V + R\,\cos\theta  }
= (c/2) \v s  + \v U + R \v s
\eeq
i.e.\ the same directional projection $R\,\v s$ enters both components. This reduces
to the pure group velocity $\dot{\v x} = (c/2)\v s + \v U$ whenever the current has no
frequency dependence ($\p_\nu \v U = 0$), e.g.\ for a barotropic current.

\subsubsection*{Other drift velocities}

Since $\v U = \v U(k)$, the material derivative of $\v U$ along a ray contains,
besides $\p_t \v U + \dot{\v x} \cdot \vn_{\v x} \v U$,  also the contribution
$\dot{\v k} \cdot \vn_{\v k} \v U$, so  $\dot \sigma$
becomes
with %$k_x = k \sin \theta$ and $k_y = k \cos \theta$ and 
$\dot{\v k} =  - (\vn \v U )\cdot \v k $
\ifdetail
\beq
\dot \sigma  &=& 2 \pi \dot \nu =   \p_t \Omega - \dot{\v U}  \cdot \v k  - \v U \cdot \dot{\v k}
\\
&=&  \p_t \Omega -  \dot U k_x + \dot V k_y  - \v U \cdot \dot{\v k}
\\ &=&  \p_t \Omega -  (\p_t U + \v u \cdot \vn_x U + \dot{\v k} \cdot \vn_{\v k} U) k_x + \dot V k_y  - \v U \cdot \dot{\v k}
\\ &=& (\p_t \v U ) \cdot \v k 
-  (\p_t \v U ) \cdot \v k -  ( \dot{\v x} \cdot \vn_x \v U) \cdot \v k 
-(\dot{\v k} \cdot \vn_{\v k} \v U) \cdot \v k
-  \v U \cdot \dot{\v k} 
\\ &=& -  ( \dot{\v x} \cdot \vn_x \v U) \cdot \v k 
-(\dot{\v k} \cdot \vn_{\v k} \v U) \cdot \v k
+  \v U \cdot  ( (\vn \v U )\cdot \v k )
\\
&=&  -   \dot{\v x} \cdot ((\vn_x \v U)  \cdot \v k) +  \v U \cdot  ( (\vn \v U )\cdot \v k )
-(\dot{\v k} \cdot \vn_{\v k} \v U) \cdot \v k
\\ &=& -( \v U  -   \dot{\v x} ) \cdot \dot \v k-(\dot{\v k} \cdot \vn_{\v k} \v U) \cdot \v k
\eeq
\fi
\beq \label{eq:omdot_Uk}
\dot \sigma &=& \p_t \Omega - \dot{\v U} \cdot \v k - \v U \cdot \dot{\v k}
 =  ( \dot{\v x}   - \v U ) \cdot \dot{\v k}
   - ( \dot{\v k} \cdot \vn_{\v k} \v U ) \cdot \v k
\eeq
\ifdetail
\beq
 ( \dot{\v k} \cdot \vn_{\v k} \v U ) \cdot \v k 
 &=&  (\dot{\v k} \cdot \vn_{\v k} U) k_x +  (\dot{\v k} \cdot \vn_{\v k}) V k_y
 \\ &=&  (\dot{k_x} \p_{k_x} U +(\dot{k_y} \p_{k_y} U ) k_x +  (\dot{\v k} \cdot \vn_{\v k}) V k_y
 \\ &=&  (\dot{k_x} \p_{k} U k_x/k +(\dot{k_y} \p_{k} U k_y/k) k_x +  (\dot{\v k} \cdot \vn_{\v k}) V k_y
 \\ &=& \p_{k} U  (\dot{\v k} \cdot \v k) k_x /k  + \p_{k} V  (\dot{\v k} \cdot \v k )  k_y/k
 =  k (\dot{\v k} \cdot \v s)  ( \p_k \v U \cdot \v s)
\eeq  
with $ \p_{k_x} k  = k_x/k$.
\fi
With $\vn_{\v k} U_i = (\p_k U_i) \, \v k /k$ and $k \, \p_k U_i = (\nu/2) \p_\nu U_i$
one finds
$( \dot{\v k} \cdot \vn_{\v k} \v U ) \cdot \v k
 = ( \dot{\v k} \cdot \v s ) \, k \lk \v s \cdot \p_k \v U \rk
 %=  ( \dot{\v k} \cdot \v s ) \,  \lk \v s \cdot (\nu/2) \p_\nu \v U \rk
 = R \, ( \v s \cdot \dot{\v k} )$
with $R$ from \eq{eq:xdot_disc}. Since also
$\dot{\v x} - \v U = (c/2 + R) \, \v s$ by \eq{eq:xdot_disc}, we obtain
\beq \label{eq:omdot_final}
\dot \sigma = 2 \pi \dot \nu = 
 (  (c/2 + R) \, \v s ) \cdot \dot{\v k}
   -  R \, ( \v s \cdot \dot{\v k} ) 
=
(c/2) \, \v s \cdot \dot{\v k}
\eeq
Finally,
\ifdetail
\beq
\tan \theta &=& k_x/k_y ~~\to ~~ \dot \theta /\cos^2\theta = (1/k_y) \dot k_x - (k_x/k_y^2) \dot k_y 
\eeq
\fi
\beq
\dot \theta = ( \cos\theta   \dot k_x -    \sin \theta \dot k_y ) k^{-1}
\eeq
results
using $\tan \theta = k_x/k_y$ and 
$\dot k_x = -  k_x \p_x U - k_y \p_x V $ and  $\dot k_y = - k_x \p_y  U   -k_y \p_y  V $.

\ifdetail
\subsubsection*{Energy transfer to mean flow}

Multiplying \eq{eq6} by $\sigma$ gives the budget of the intrinsic energy
$E = \sigma A$. Since $\sigma = 2 \pi \nu$ is the spectral coordinate, it
commutes with $\p_t$, $\vn_{\v x}$ and $\p_\theta$; only the frequency
advection generates
$\sigma \, \p_\nu ( \dot \nu A) = \p_\nu ( \dot \nu E) - \dot \sigma A$
so that
\beq \label{eq:ebudget}
\pt{E} + \vn_{\v x} \cdot ( \dot{\v x} E ) + \p_\theta ( \dot \theta E )
 + \p_\nu ( \dot \nu E )
 &=& \sigma \dot A + \dot \sigma A
\eeq
All terms on the left are conservative transport of intrinsic energy; on the
right, $\sigma \dot A$ is the energy of the source terms and $\dot \sigma A$
is the local energy exchange with the mean flow mediated by propagation.
Multiplying \eq{eq6} instead by $\Omega(\v x, \theta, f, t)$
yields 
\ifdetail
\beq 
  \Omega  \pt{ A}    +  \Omega  \vn_{\v x}  \cdot   (\dot{\v x} A ) 
  + \Omega \p_{\theta}  (\dot{\theta} A  )
 + \Omega \p_{\nu} ( \dot{\nu} A)
 &=& \Omega  \dot A 
 \\
  \pt{ \Omega   A}  - A \p_t \Omega 
  +  \vn_{\v x}  \cdot ( \Omega  \dot{\v x} A )  
  + \p_{\theta}  (\Omega  \dot{\theta} A  )
  +\p_{\nu} (  \Omega  \dot{\nu} A) 
  - A ( \dot{\v x}    \cdot \vn_x  \Omega 
   + \dot{\theta} \p_{\theta} \Omega  
  +  \dot{\nu}  \p_\nu  \Omega )
   &=& \Omega  \dot A 
   \eeq
   \beq
  \pt{ \Omega   A} 
  +  \vn_{\v x}  \cdot ( \Omega  \dot{\v x} A )  
  + \p_{\theta}  (\Omega  \dot{\theta} A  )
  +\p_{\nu} (  \Omega  \dot{\nu} A) 
    &=& \Omega  \dot A  + A \p_t \Omega 
\eeq
\fi
\beq \label{eq:absE}
 \p_t   E + \vn_{\v x} \cdot ( \dot{\v x} \Omega A )
 + \p_\theta ( \dot \theta \Omega A ) + \p_\nu ( \dot \nu \Omega A )
 &=& \Omega \dot A  - \v U \cdot \v k \p_t A
\eeq
using $\dot{\v x} \cdot \vn_{\v x} \Omega + \dot \theta \, \p_\theta \Omega
 + \dot \nu \, \p_\nu \Omega = 0$.

Except for flux divergencies the
exchange $C$ thus proceeds through three routes: (i) the refraction transfer
$ \dot \sigma A$ to the ocean,
carried by the frequency advection,
(ii) the source-mediated part
$\v U \cdot \v k \, ( S_{in} + S_{nl} ) $,
handed over as $\p_t \v u_S$-work, and (iii) the dissipation part routed
through the stress $\v \tau$ related to $\v U \cdot \v k \,  S_{out}$.
In a horizontally homogeneous column all propagation terms vanish, $\p_t A = \dot
A$, and $C$ collapses to routes (ii) and (iii): energy flows to the ocean
with $\dot \sigma \equiv 0$, since each action quantum created or destroyed
by the sources carries the mean-flow energy $\v U \cdot \v k$ alongside its
intrinsic energy $\sigma$.


\fi
\subsection*{High frequency cut-off}

Since the growth rates $\gamma_{in}$ and $\gamma_{out}$ increase strongly with frequency,
which would require a very small time step at large $\nu$,
the wave model is only calculated up to  a cut-off frequency $\nu_c$,
and for the higher frequencies $\nu > \nu_c$ 
a Phillips spectrum with a $\nu^{-5}$ tail of the wave
spectrum is  applied, the so-called high-frequency tail. It has the form 
$E_{hi}(\nu)  =   E(\nu_c,\theta) \nu_c^5 \nu^{-5}$. 
The cut-off frequency $\nu_c$ is calculated as in WAMDI (1988) 
\beq
\nu_c = \max( 2.5 \bar \nu, 3 \nu_{pm})
\eeq
with the mean frequency 
$2 \pi \bar \nu =  \int E d \theta d\nu/ \int (E/\sigma) d \theta d\nu$, the ''wind sea'' frequency $2 \pi \nu_{pm} =  g/(28 u_{\star})$,
and the friction velocity $u_\star$.
For $\nu > \nu_c$ it is assumed that \eq{eq6} becomes local and steady, such that the
source terms balance locally in space and time and in $\nu$ and $\theta$, i.e.\ $S |_{\nu>\nu_c}=0$.
The per-column prognostic cut-off and the imposed $\nu^{-5}$ tail are activated by the
attribute \verb|mask_energy| (default off); the resulting prognostic/diagnostic bin mask
is also the basis of the equilibrium tail-flux booking of \verb|tail_flux| $=$ \verb|'input'|
described below, which therefore requires \verb|mask_energy| to be active.


\subsection*{Source terms}

The source terms $S_{in}$ and $S_{out}$ are parameterized by growth rates,
$ S_{in} = \gamma_{in} A$ and $S_{out}=  \gamma_{out} A$, 
for which the simplest forms  following  WAMDI (1988) are used
\beq \label{eq:gammas}
 \gamma_{out} = -  \beta_{o} g^{-4}  \bar E^2 \bar \sigma^7 \, \,     \sigma^2  ~&,&~
 \gamma_{in} =   \sigma\frac{\rho_a}{\rho_o} \times \, 
    0.25 \max\left( 0,
     28 (u^*/c)  \cos(\theta - \theta*) -1  \right) )
 \eeq
with mean energy $\bar E = \int E d \theta d\nu$,  mean frequency
$\bar \sigma = \bar E/ \int (E/\sigma) d \theta d\nu$,
air,water density  $\rho_a,\rho_o$, atmospheric friction velocity magnitude $u^{\star}$,
with direction $ \theta^*$, deep water phase velocity $c = g/  \sigma$,
and the parameter $\beta_{o} = 2.56 $.
Other parameterisations are implemented as options.
The mean frequency is clipped to the resolved spectral range,
\beq
\bar\sigma \leftarrow \min\!\left(\max(\bar\sigma,\, \sigma_0),\, \sigma_{N_\nu-1}\right)
\eeq
The non-linear source term $S_{nl}$ is the discrete interaction approximation (DIA)
of Hasselmann et al.\ (1985), taken from WAM (cycle 7).
A fortran version directly from the original WAM legacy code, 
and two python versions (one with jax/jit acceleration, one pure numpy) are available.
The attribute ''snonlin\_backend'' determines which one to use ('fortran'/'jax'/'numpy'), default is 'fortran';
all versions give identical results.


\subsection*{Stokes velocity}

Stokes velocity is  
\beq
\v u_S(z) =  \int \int 2 (k/\sigma) e^{2kz} \v k   E(\v k)  d \v k  
= (2/g) \int_0^\infty  \sigma^3  e^{2k z } \int_0^{2 \pi}   \v s E  d\theta d \nu 
\eeq
with direction vector $\v s=(\sin \theta,\cos \theta)$. The frequency integral is calculated over the
prognostic range of $\nu$, and the contribution by the high frequency tail is added.
The latter part is given by
  \beq \label{eq:ushi}
  \v u_S^{hi} = 
  (2/g) \int_{\nu_c}^\infty \int_0^{2 \pi}    \sigma^3  e^{2z \sigma^2/g}    \v s  E(\nu_c,\theta)  \nu_c^5  \nu^{-5} d\theta d \nu 
  \approx  (2/g)  \nu_c  \int_0^{2 \pi}   \sigma_c^3 \v s  E(\nu_c,\theta)  d\theta
%  \\  \approx 
 %   (2/g) \int_{\nu_c}^\infty \sigma^3 \nu_c^5  \nu^{-5} d \nu  \int_0^{2 \pi}   \v s  E(\nu_c,\theta) d\theta 
 %   =   (2/g)  (2 \pi)^3\nu_c^5 /\nu_c  \int_0^{2 \pi}   \v s  E(\nu_c,\theta)  d\theta
 \eeq 
 assuming that the uppermost model level thickness consumes the factor $ e^{2z \sigma^2/g} $
 and with $\sigma_c = 2 \pi \nu_c$. Below the uppermost vertical level, $  \v u_S^{hi} =0$ in this
 approximation.
 For small effective vertical decay scale of the tail contribution or large
 vertical levl thickness, this approximation
 gets worse.
The frequency integral in \eq{eq:ushi} can in fact be carried out
exactly, retaining the full depth dependence of the exponential. With
$\sigma_c = 2\pi \nu_c$ the result is
%
\begin{equation}
  u_S^{hi}(z)
  = \frac{2}{g}\, \nu_c \int_0^{2\pi}\! \mathbf{s}\,E(\nu_c,\theta)\,
    \sigma_c^{3}\left[\,
      e^{2z\sigma_c^{2}/g}
      - \sigma_c\sqrt{-\frac{2\pi z}{g}}\;
        \erfc\!\left(\sigma_c\sqrt{-\frac{2z}{g}}\right)
    \right] d\theta ,
  \label{eq:ushi-exact}
\end{equation}
%
valid for $z\le 0$. The approximation \eq{eq:ushi} is recovered in the
near-surface limit $\sigma_c^{2}|z|/g \ll 1$, where the $\erfc$ term
vanishes and $e^{2z\sigma_c^{2}/g}\to 1$, leaving the bracket equal to
$\sigma_c^{3}$ and hence the prefactor $\frac{2}{g}\nu_c\,\sigma_c^{3}$ of
\eq{eq:ushi}. Note that WAM (cycle 7) uses the approximated version;
the difference grows with the ratio of the uppermost level thickness to the
vertical decay scale of the tail.
 
\ifdetail
 
%
Starting from the high-frequency tail of \eq{eq:ushi} with the Phillips
form $E_{hi}(f)=E(\nu_c,\theta)\,\nu_c^{5}\nu^{-5}$, the per-direction frequency
integral to evaluate is
%
\beq
  I(z) \;=\; \int_{\nu_c}^{\infty}\sigma^{3}\,
              e^{2z\sigma^{2}/g}\,\nu_c^{5}\,\nu^{-5}\,d\nu ,
  \qquad \sigma = 2\pi \nu .
\eeq
%
Substituting $\sigma^{3}=(2\pi)^{3}\nu^{3}$ and writing
$b \equiv 8\pi^{2}z/g \le 0$ (since $z\le 0$, so $e^{b\nu^{2}}$ decays),
%
\beq
  I(z) \;=\; (2\pi)^{3} \nu_c^{5}\int_{\nu_c}^{\infty} \nu^{-2}\,e^{b \nu^{2}}\,d\nu .
\eeq
%
The remaining integral is evaluated by parts, differentiating $e^{b\nu^2}$
and integrating $\nu^{-2}$:
%
\beq
  \int_{\nu_c}^{\infty} \nu^{-2} e^{b \nu^{2}}\,d\nu
  = \Bigl[-\nu^{-1} e^{b \nu^{2}}\Bigr]_{\nu_c}^{\infty}
    + 2b\int_{\nu_c}^{\infty} e^{b \nu^{2}}\,d\nu
  = \frac{e^{b \nu_c^{2}}}{\nu_c}
    - \sqrt{-\pi b}\;\erfc\!\left(\nu_c\sqrt{-b}\right),
\eeq
%
where the boundary term at infinity vanishes ($b<0$) and the Gaussian
integral gives
$\int_{\nu_c}^{\infty} e^{b \nu^{2}}\,d\nu
   = \frac{1}{2}\sqrt{-\pi/b}\,\erfc(\nu_c\sqrt{-b})$.
Hence
%
\beq
  I(z) = (2\pi)^{3} \nu_c^{5}
         \left[\frac{e^{b \nu_c^{2}}}{\nu_c}
               - \sqrt{-\pi b}\;\erfc\!\left(\nu_c\sqrt{-b}\right)\right].
\eeq
%
Re-expressing in terms of $\sigma_c = 2\pi \nu_c$ using
$b \nu_c^{2} = 2z\sigma_c^{2}/g$,
$\;\nu_c\sqrt{-b} = \sigma_c\sqrt{-2z/g}\,$, and
$\;(2\pi)^{3}\nu_c^{5}\sqrt{-\pi b}
   = \sigma_c^{3}\nu_c\cdot\sigma_c\sqrt{-2\pi z/g}\,$ gives
%
\beq
  I(z) = \sigma_c^{3} \nu_c
         \left[\,e^{2z\sigma_c^{2}/g}
               - \sigma_c\sqrt{-\frac{2\pi z}{g}}\;
                 \erfc\!\left(\sigma_c\sqrt{-\frac{2z}{g}}\right)\right].
\eeq
%
Inserting $I(z)$ back, together with the prefactor $2/g$ and the angular
integral $\int_0^{2\pi}\mathbf{s}\,E(\nu_c,\theta)\,d\theta$, yields
\eq{eq:ushi-exact}. As $z\to 0^{-}$ the exponential tends to unity and
the $\erfc$ argument to zero with its $\sqrt{-z}$ prefactor, so
$I\to\sigma_c^{3}\nu_c$, recovering \eq{eq:ushi}.
 \fi
 
\subsection*{Wave stress}

The stress force in the ocean is given, with the dissipation energy source
$\dot E_{out} = \gamma_{out} E$ in $(\theta,\nu)$ coordinates, by
\beq \label{eq:dztau}
 \p_z \v \tau %= - \int Q \v k S_{out} d \v k 
 =  - (2 /g) \int_0^\infty  e^{2kz}  \sigma^3 \int_0^{2 \pi}  \v s  \dot E_{out}  d \theta d\nu 
  \eeq
The stress itself is 
\beq \label{eq:tauo}
 \v \tau_o = %- \int  \v k S_{out} d \v k = 
  -   \int_0^\infty\sigma   \int_0^{2 \pi}     \v s \dot E_{out}  d \theta d\nu
   \eeq 
The contribution  by the high-frequency tail assuming the WAMDI dissipation scale is
\beq \label{eq:dztau-hi}
 \p_z \v \tau^{hi}  
 %=  - (2 /g) \int_{\nu_c}^\infty  e^{2kz}  \sigma^3  \gamma_{out} \int_0^{2 \pi}  \v s  E  d \theta d\nu \\ 
 %=  2 \nu_c^5  \beta_{o} g^{-5}  \bar E^2 \bar \sigma^7 (2 \pi)^5 \int_{\nu_c}^\infty  e^{2kz}   d\nu
% \int_0^{2 \pi}  \v s   E(\nu_c,\theta)   d \theta \\
 =  2 \nu_c^5  \beta_{o} g^{-5}  \bar E^2 \bar \sigma^7 (2 \pi)^5 \int_{\nu_c}^\infty  e^{ (8 \pi^2/g) z  \,  \nu^2 }   d\nu
 \int_0^{2 \pi}  \v s   E(\nu_c,\theta)   d \theta 
  \eeq
   with $E_{hi}(\nu)  =   E(\nu_c,\theta) \nu_c^5 \nu^{-5}$ and 
   $\gamma_{out} = -  \beta_{o} g^{-4}  \bar E^2 \bar \sigma^7 \sigma^2 $ and the finite integral
   \beq \label{eq:erfint}
    \int_a^{\infty \:}e^{-bx^2}dx = \frac{ \sqrt{\pi }- \sqrt{\pi }\mbox{erf}\left(a\sqrt{b}\right)} { 2\sqrt{b}}
   \eeq
  For the stress we find for the high-frequency tail 
   \beq \label{eq:tauo-hi}
 \v \tau_o^{hi}  =    
 %  -   \nu_c^5 \int_{\nu_c}^\infty\sigma  \gamma_{out}  \nu^{-5} d\nu \int_0^{2 \pi}     \v s  E(\nu_c,\theta)d \theta 
 %\\ = 
  (2 \pi)^3 \nu_c^4 \beta_{o} g^{-4}  \bar E^2 \bar \sigma^7 
 \int_0^{2 \pi}     \v s   E(\nu_c,\theta)  d \theta 
   \eeq 

An alternative formulation of the tail stress follows from the equilibrium
assumption that underlies the imposed $\nu^{-5}$ shape itself: for $\nu > \nu_c$
the sources balance, $\dot E_{in} + \dot E_{nl} + \dot E_{out} = 0$, so the
flux to the ocean, $-\dot E_{out}$, can be evaluated as
$\dot E_{in} + \dot E_{nl}$ instead of extending the WAMDI dissipation scale
(which is calibrated on the resolved spectrum and underestimates the
dissipation that maintains the equilibrium range). The input part is
\beq \label{eq:tauo-hi-in}
 \v \tau_o^{hi,in} = \int_{\nu_c}^\infty \sigma\, \gamma_{in}\, \nu_c^5 \nu^{-5}
   \, d\nu \int_0^{2\pi} \v s\, E(\nu_c,\theta)\, d\theta
\eeq
with $\gamma_{in}$ from \eq{eq:gammas}. Writing $28 (u^\star/c)
\cos(\theta - \theta^*) = a \nu$ with $a(\theta) = 28 u^\star 2\pi
\cos(\theta-\theta^*)/g$, the input is active only above
$\nu_1 = \max(\nu_c, 1/a)$, and the frequency integral is elementary,
$\int_{\nu_1}^\infty (a \nu^{-2} - \nu^{-3})\, d\nu = a/\nu_1 - 1/(2 \nu_1^2)$, giving
per direction
\beq \label{eq:tauo-hi-in2}
 \v \tau_o^{hi,in} = \frac{\rho_a}{\rho_o} \frac{(2\pi)^2 \nu_c^3}{4}
 \int_0^{2\pi} \v s\, E(\nu_c,\theta)
 \left\{ \begin{array}{ll} a \nu_c - \frac{1}{2}, & a \nu_c \ge 1\\[2pt]
 \frac{1}{2}\, (a \nu_c)^2, & 0 < a \nu_c < 1 \end{array} \right\}  d\theta
\eeq
(continuous at $a \nu_c = 1$; zero for $\cos(\theta-\theta^*) \le 0$). The
integrand decays like $\nu^{-2}$, i.e.\ the integral converges exactly as
benignly as the WAMDI form in \eq{eq:tauo-hi}; note that the corresponding
integral with the dissipation of Janssen (1994; see also Komen et al.\ 1994,
and the ECMWF IFS documentation, Part VII) diverges strongly because of its
$(k/\bar k)^4$ term -- its $(k/\bar k)^2$ term behaves like the WAMDI form --
which is the reason the WAMDI scale is used in \eq{eq:tauo-hi}. The nonlinear part uses the conservation of the exact
transfer over the full spectrum, $\int_0^\infty \sigma \int \v s\,
\dot E_{nl}\, d\theta\, d\nu = 0$, so that
\beq \label{eq:tauo-hi-nl}
 \v \tau_o^{hi,nl} = \int_{\nu_c}^\infty \sigma \int_0^{2\pi} \v s\,
 \dot E_{nl}\, d\theta\, d\nu
 = - \int_0^{\nu_c} \sigma \int_0^{2\pi} \v s\, \dot E_{nl}\, d\theta\, d\nu
\eeq
i.e.\ the cascade flux into the tail equals minus the resolved-range
integral of the computed nonlinear transfer (the ``cascade credit''). For
completeness, the resulting total stress on the ocean in this
formulation is
\beq \label{eq:tauo-in-tot}
 \v \tau_o = - \int_0^{\nu_c} \sigma \int_0^{2\pi} \v s
 \left( \dot E_{out} + \dot E_{nl} \right) d\theta\, d\nu
 \; + \; \v \tau_o^{hi,in}
\eeq
where the first term is the resolved dissipation of \eq{eq:tauo} restricted
to the prognostic range with the cascade credit \eq{eq:tauo-hi-nl} folded
in, and the second term is the analytic tail input \eq{eq:tauo-hi-in2}.

The depth-distributed version (the ``stress work'' of \eq{eq:dztau}) follows
with the same replacement of the tail rate. With
$-\dot E_{out} = \gamma_{in} E(\nu_c,\theta) \nu_c^5 \nu^{-5}$ above $\nu_c$ and
$\sigma^3 \gamma_{in} = \frac{\rho_a}{\rho_o} \frac{(2\pi)^4}{4}
\nu^4 (a \nu - 1)$,
\beq \label{eq:dztau-hi-in}
 \p_z \v \tau^{hi,in} = \frac{2}{g} \frac{\rho_a}{\rho_o}
 \frac{(2\pi)^4}{4} \nu_c^5 \int_0^{2\pi} \v s\, E(\nu_c,\theta)
 \int_{\nu_1}^\infty e^{ (8\pi^2/g) z\, \nu^2 }
 \left( a - \nu^{-1} \right) d\nu\, d\theta
\eeq
with $\nu_1 = \max(\nu_c, 1/a)$ as before. The $a$--term is the finite integral
\eq{eq:erfint}; the second term requires the companion
\beq \label{eq:e1int}
 \int_{a_0}^{\infty} x^{-1} e^{-b x^2} dx = \frac{1}{2} E_1( a_0^2\, b )
\eeq
with the exponential integral $E_1$. The vertical integral of
\eq{eq:dztau-hi-in} recovers \eq{eq:tauo-hi-in2} exactly. The cascade
contribution \eq{eq:tauo-hi-nl} is deposited at the surface: it is
dissipated by tail waves with $\nu \ge \nu_c$, whose deposition scales
$g/2\sigma^2$ are centimetres, far below any ocean grid spacing. Both
formulations are implemented, selected by the switch \verb|tail_flux|
$\in$ \{\verb|'wamdi'|, \verb|'input'|\}; the discrete treatment, including
the diagnostic bins inside the grid, is described below.

\subsubsection*{Discretisation of the wave--ocean coupling and the energy exchange}

All three coupling quantities of \eq{eq_Q} -- the Stokes velocity $\v u_S$,
the stress force $\p_z \v \tau$ and the Doppler velocity $\v U$ -- involve the
same vertical kernel $Q(k,z) = 2 k e^{2kz}$ with $\int_{-\infty}^0 Q \, dz =
1$: $\v u_S$ and $\p_z \v \tau$ assign spectral quantities to a depth profile,
while $\v U$ projects the Lagrangian velocity back onto the spectral grid. The
exchange term $C$ of \eq{eq_C} closes because $\v U$ is the adjoint of the
profile assignment, i.e.\ because the same $Q$ appears on both sides of the
depth integral. The discretisation preserves this structure: one common
discrete kernel is used in all three methods of
\verb|./pyOM3/wim/interface.py| (\verb|calc_stokes|, \verb|calc_stress_work|,
\verb|calc_doppler_vel|), given by the cell integral of $Q$ over each ocean
cell of thickness $\Delta z$ with centres $z_k$ and faces $z_{k \pm 1/2} = z_k
\pm \Delta z/2$ (uppermost face at $z=0$),
\beq \label{eq:Qdz}
 \bar Q_k \Delta z = \int_{z_{k-1/2}}^{z_{k+1/2}} Q \, dz 
 = e^{2 k z_{k+1/2}} - e^{2 k z_{k-1/2}} 
\eeq
in place of the point value $Q(z_k) \Delta z$. The cell integral is essential:
the decay scale $g/(2\sigma^2)$ of the kernel is of the order of centimetres
for the high frequencies of the prognostic range and for the tail, far below
any ocean level thickness, so that point sampling loses a large part of the
coupling amplitude (about half of the stress force for $\Delta z = 1 \, m$ and
a developing wind sea). The weights \eq{eq:Qdz} telescope over the column,
$\sum_k \bar Q_k \Delta z = 1 - e^{-2 k D}$ for column depth $D$, with three
consequences that hold exactly for any $\Delta z$: the column-integrated
stress force equals the surface stress $\v \tau_o$ of \eq{eq:tauo}, the
discrete Stokes transport $\sum_k \v u_S(z_k) \Delta z$ equals its analytic
value, and a depth-independent (barotropic) current is projected onto the
Doppler velocity without attenuation, $\v U = \v u_L (1 - e^{-2kD})$. The
frequency derivative $\p_\nu \v U$ entering the radial refraction term
\eq{eq:xdot_disc} is taken as the analytic derivative of the discrete kernel,
$\p_\nu (\bar Q_k \Delta z) = (2k/\nu) \, 2 ( z_{k+1/2} e^{2 k z_{k+1/2}} -
z_{k-1/2} e^{2 k z_{k-1/2}} )$, so that refraction sees the same Doppler
projection as the advection.

For the high-frequency tails the same principle is applied to the exact tail
profiles. The tail of the stress force \eq{eq:dztau-hi} contains the frequency
integral $W(t) = \int_{\nu_c}^\infty e^{-b t \nu^2} d\nu$ with $t = -z \ge 0$ and $b
= 8 \pi^2/g$, evaluated with \eq{eq:erfint}; its antiderivative is
\beq \label{eq:tailH}
 H(t) = \int_0^t W(t') dt' = \frac{\sqrt{\pi}}{b} \left[ \sqrt{b t} \, 
 \mbox{erfc} \left( \nu_c \sqrt{b t} \right) 
 + \frac{1 - e^{-\nu_c^2 b t}}{\nu_c \sqrt{\pi}} \right]
\eeq
and each cell receives $\left( H(-z_{k-1/2}) - H(-z_{k+1/2}) \right)/\Delta z$
in place of the point value. Since $H(\infty) = 1/(b \nu_c)$, the column sum of
the tail force equals its exact depth integral $\v \tau_o^{hi}$ of
\eq{eq:tauo-hi}, using $2 (2 \pi)^5 \nu_c^5 g^{-5} H(\infty) = (2 \pi)^3 \nu_c^4
g^{-4}$; the prefactor $(2 \pi)^5$ in \eq{eq:dztau-hi} is essential for this
consistency. The Stokes tail \eq{eq:ushi-exact} is treated the same way: with
$a = 2 \sigma_c^2/g$ its depth profile $B(t) = \nu_c [ e^{-a t} - \sqrt{\pi a t}
\, \mbox{erfc}(\sqrt{a t}) ]$ has the antiderivative ($S = \sqrt{t}$)
\beq \label{eq:Istokes}
 I(t) = \nu_c \left[ \frac{1 - e^{-a t}}{a} - 2 \sqrt{\pi a} \left( 
 \frac{S^3}{3} \mbox{erfc}(\sqrt{a} S) 
 + \frac{1 - (a S^2 + 1) e^{-a S^2}}{3 \sqrt{\pi} \, a^{3/2}} \right) \right]
\eeq
with $I(\infty) = \nu_c/(3 a)$, so that the discrete tail transport equals the
exact value $(2\pi/3) \nu_c^2 \int \v s E(\nu_c,\theta) d\theta$. In tests the
column sums reproduce $\v \tau_o$ and the
Stokes transport to machine precision; for a shallow column below a
developed wind sea, however, a substantial part of the retained fraction $1 - e^{-2kD}$
of the wave momentum flux can be deposited below the column bottom. For a
closed momentum budget the attribute \verb|closed_bottom| extends the
integration of the bottom cell to $z \to -\infty$, in the shared kernel
($e^{2 k z_{bot}} \to 0$ in \eq{eq:Qdz}) as well as in both tail
antiderivatives ($H, I \to H(\infty), I(\infty)$), so that $\sum_k \bar Q_k
\Delta z = 1$ exactly: the bottom cell absorbs the deep part of all coupling
profiles, no wave momentum is lost, and the barotropic Doppler projection
becomes exact, $\v U = \v u_L$, at all frequencies (with the open bottom the
low-frequency bins of a shallow column see an almost vanishing Doppler
current). Since the kernel remains shared, the adjoint identity
\eq{eq:Cdisc} below is unaffected. The closure amounts to assuming $\v u_L$
constant below the column floor and is the recommended setting for shallow
test columns; for deep columns ($2 k D \gg 1$) it makes no difference.

The energy exchanged between ocean and waves is controlled by the same
kernel. On the ocean side, with the adjoint pair \eq{eq:Qdz} the discrete
identity
\beq \label{eq:Cdisc}
 \sum_k \v u^L_h(z_k) \cdot \p_t \v u_S(z_k) \, \Delta z 
 = \int \v U \cdot \v k \; \p_t A \; d \v k
\eeq
holds exactly (to machine precision) for any spectral tendency $\p_t A$ and
any $\Delta z$, over the prognostic frequency range (the tails are diagnostic
and remain outside this closure). The ocean therefore receives, per coupling
step, exactly the energy $\int \v U \cdot \v k \, \p_t A \, d \v k$
corresponding to the actual discrete evolution of the wave field, and the
handoff itself introduces no discretisation error.

On the wave side, the counterpart of this energy must be extracted from (or
added to) the intrinsic wave energy $\int \sigma A \, d \v k$ by the
propagation operator. For steady $\v U$ the rays conserve $\Omega = \sigma +
\v U \cdot \v k$, so the continuous propagation operator satisfies $\int
\Omega \, \p_t A |_{prop} \, d \v k = 0$ and the intrinsic energy change $\int
\sigma \, \p_t A$ balances the ocean gain $-\int \v U \cdot \v k \, \p_t A$
identically. Discretely, the flux-form advection conserves the action $\int A
\, d\v k$ to machine precision, but not $\int \Omega A \, d \v k$: the
nonlinear Superbee limiter leaves a residual (the frequency advection velocity
$\dot \nu$ is already located on the $u$-centred frequency faces, so face
staggering is not the cause). In tests with a wave packet refracted across a jet
(moderate Doppler shift, $\max |\v U \cdot \v k|/\sigma \approx 0.3$) this
$\Omega$-conservation defect amounts to about $1 \, \%$ of the gross local
transfer rate $\int |\dot \sigma A| \, d\v k$ for the Superbee scheme,
consistent with the $\mathcal{O}(0.5 \, \%)$ energy overshoot observed for such a
packet. First-order upwind advection is far worse (about
$12 \, \%$, since its frequency diffusion moves action across the $\sigma$
gradient), an unlimited second-order flux reduces the defect to about $0.2 \,
\%$. The $x$-sweep is exactly energy-neutral for a zonal jet; the defect
resides in the $y$, $\nu$ and $\theta$ sweeps.

The residual is removed exactly by a conservative energy fixer
(attribute \verb|energy_consistent| in \verb|propagation.py|). Writing the
discrete advection in flux form per phase-space grid cell $i$ as
$\pt{(A_i \Delta_i)} = -\sum_{faces} \delta F_i$, with the phase-space
volume $\Delta_i$ of the cell and $F$ the total (limited) action transport
through a cell face in $\v x$, $\nu$ or $\theta$, the transports cancel
pairwise in the sum over all cells and action is conserved to machine
precision for any advection scheme. Multiplying by $\Omega_i$, summing over
all grid cells and shifting the sums equivalent to a discrete partial
integration yields
\beq \label{eq:edefect}
 \pt{} \sum_i \Omega_i A_i \Delta_i = \sum_i \sum_{faces} F_i \, \delta \Omega_i
 \equiv D
\eeq
with the forward difference $\delta \Omega$ of $\Omega$ across each cell
face. Boundary terms vanish for closed and periodic boundaries (not for
reflective ones, for which no correction is made). The possible leakage $D$
is split into contributions $\sum_{j,k} D_{j,k} = D$ per water column at
discrete position $(j,k)$, with the shared horizontal face terms split
half-half between neighbouring columns, and each $D_{j,k}$ is cancelled by a
small corrective action transport $\tilde F_m$ across the frequency faces of
that column,
\beq \label{eq:efix}
 \tilde F_m = \lambda \, \bar A_m \, \delta_\nu \Omega_m ~~,~~
 \lambda(\v x) = - \frac{ D_{j,k} }
                       { \sum_m \bar A_m \, (\delta_\nu \Omega_m)^2 }
\eeq
with $m$ running over all $\nu$ faces and directions $\theta$ at position
$(j,k)$, and the weight $\bar A_m$ given by the face average of the action
content of the two adjacent cells. Being a transport, $\tilde F$ leaves the
action conservation untouched, but by the same partial integration it changes
the energy by $\sum_m \tilde F_m \, \delta_\nu \Omega_m = -D_{j,k}$. The
frequency direction is chosen for the corrective flux since it is the only
one in which $\delta \Omega$ does not vanish without currents
($\delta_\nu \Omega \approx \delta_\nu \sigma > 0$ for small Doppler
shift), and it is where the leakage is generated by action crossing the
frequency gradient of $\Omega$. The factor $\bar A_m$ ties the correction to
the level of action, so that it acts as a small correction of the frequency
propagation velocity rather than as a source term and positivity is not
threatened, and the factor $\delta_\nu \Omega_m$ makes the denominator
sign-definite, so that $\lambda$ exists whenever action is present at that
position. Minimising the action-weighted amplitude
$\sum_m \tilde{F}_m^2 / \bar A_m$ subject to this constraint yields the same
form by the standard Lagrange-multiplier condition: among all conservative
frequency-flux corrections that restore the energy, \eq{eq:efix} is the
smallest one in the action-weighted least-squares sense, i.e.\ the correction
perturbs the solution as little as possible.

The correction is local (one scalar per position $\v x$, no global
communication) and dormant without currents: for $\v U = 0$ we have
$\Omega = \sigma$, no frequency transports exist, and $D_{j,k} = 0$ so that
$\lambda = 0$. In tests the correction amounts to about $1 \, \%$ of the
frequency transports, the tendency-level defect drops from $0.9 \, \%$ of the
gross transfer to machine precision, and in long integrations of a wave
packet through a jet the energy $\int \Omega A \, d \v k$ is conserved
exactly (Adams--Bashforth time stepping preserves this, since each stored
tendency is corrected individually and $\Omega$ is constant while $\v U$ is
frozen between coupling updates), whereas without the fixer the spurious
drift amounts to roughly half of the genuine exchange. The transfer from
ocean to waves is thus exact on both sides; a structurally energy-conserving
advection would achieve the same without an a posteriori correction and
remains an option for future work.
   
 \subsection*{The coupling to the ocean}
 
 
 Stokes velocity, wave stress, and wave stress force are calculated in \verb|./pyOM3/wim/interface.py|.
 The methods there may be called and results passed to the ocean model.
 The velocity variable of the coupled ocean model is interpreted as the
 three-dimensional Lagrangian mean velocity $\v u_L = \v u + \v u_S$, with the
 Eulerian mean velocity $\v u$ and the Stokes velocity $\v u_S$, which is the most
 natural choice for energetic consistency (it also follows the example of interpreting
 the model's velocity as the isopycnal thickness weighted velocity to account for the
 effect of meso-scale eddies). Starting from the Craik--Leibovich equations of motion,
 the momentum equation for $\v u_L$ becomes
\beq \label{eqM}
\p_t \v u_L  -  \v u_L \times \vn \times (\v u_L-\v u_S)
 + f \v z \times \v u_L  = - \vn \pi + \p_t \v u_S + (\p_z \v  \tau,0)
\eeq
with the three-dimensional operator $\vn = (\vn_{\v x}, \p_z)$ and a gradient force
$\vn \pi$ collecting the pressure gradient and wave-driven gradient components.
The ocean model thus receives the wave-driven force $\p_t \v u_S$, the vortex force
\beq \label{eq:vf_cont0}
\v F = - \v u_L \times \vn \times \v u_S
\eeq
with $\v u_L \cdot \v F = 0$, and the wave momentum flux (stress), either as the
surface stress $\v \tau_o$ into the uppermost grid cell or as the interior stress
force $\p_z \v \tau$; the tracers are advected by $\v u_L$.
Note that the Coriolis term acts on the full $\v u_L$: with $w = 0$, identical
results are obtained using $\v u$ as model velocity and adding the Stokes--Coriolis
force $-f \v z \times \v u_S$ instead of $\p_t \v u_S$, but with $w \neq 0$ the
Lagrangian formulation is required to stay energetically consistent. The scalar
product of \eq{eqM} with $\v u_L$ contains no transfer terms related to the vortex
force and the Coriolis--Stokes force; the direct transfer to or from the waves is by
$\v u_L \cdot \p_t \v u_S$ and by the stress work $\v u_h^L \cdot \p_z \v \tau$,
which together form the exchange $C$ of \eq{eq_C}.

\subsubsection*{Discrete vortex force for horizontally uniform $\v u_S$ (nh\_simple)}

For a horizontally uniform Stokes velocity $\v u_S = (u_S(z),v_S(z),0)$, as used in
the non-hydrostatic model nh\_simple, the wave-driven force of \eq{eq:vf_cont0} becomes
\beq
\v F =  - \v u_L \times \vn \times \v u_S =
 \Vec{ \; w \, \pz{u_S} \\  \;w \, \pz{v_S} \\ -  u \,\pz{u_S} - v\, \pz{v_S} }
 ~~,~~\v u_L \cdot\v F = 0
 \label{eq:vf_cont}
\eeq
writing $\v u_L = (u,v,w)$ in this section for brevity.
On the C-grid, $u$, $v$ and $w$ are staggered at the cell faces
$(i+\frac{1}{2},j,k)$, $(i,j+\frac{1}{2},k)$ and $(i,j,k+\frac{1}{2})$,
and the vertical Stokes shear is placed on the $w$-levels,
\beq
S^u_{k+\frac{1}{2}} = \frac{ u_{S,k+1}-u_{S,k} }{\Delta z} ~~,~~
S^v_{k+\frac{1}{2}} = \frac{ v_{S,k+1}-v_{S,k} }{\Delta z}
\eeq
with $S^u = S^v = 0$ at the rigid lid and the flat bottom.
Introducing the two-point averages
$\Bar{ \phi }^{x}_{i+\frac{1}{2}} = ( \phi_{i} + \phi_{i+1} )/2$
and analogous in $y$ and $z$, the discrete force $\v F=(F^u,F^v,F^w)$ is
\beq
F^u_{i+\frac{1}{2},j,k} = \Bar{ \; \Bar{w}^{x} \, S^u \; }^{z} ~~,~~
F^v_{i,j+\frac{1}{2},k} = \Bar{ \; \Bar{w}^{y} \, S^v \; }^{z}
~~,~~
F^w_{i,j,k+\frac{1}{2}} =
 - \Bar{ \; \Bar{u}^{x} \; }^{z} S^u  - \Bar{ \; \Bar{v}^{y} \; }^{z} S^v
\label{eq:vf_disc}
\eeq
i.e.\ the products $\Bar{w}^{x} S^u$ are formed on the cell edges
$(i+\frac{1}{2},j,k+\frac{1}{2})$ and averaged back to the $u$-points,
while for $F^w$ the horizontal velocities are averaged first to the cell
centers and then to the $w$-levels where the shear is located.
The averaging operators in \eq{eq:vf_disc} are pairwise transposed:
each average that moves $w$ towards $u$ in $F^u$ is mirrored by an average
moving $u$ towards $w$ in $F^w$. With the summation-by-parts identity for
the two-point average
\beq
\sum_i \psi_{i+\frac{1}{2}} \, \Bar{\phi}^{x}_{i+\frac{1}{2}}
= \sum_i \phi_{i} \, \Bar{\psi}^{x}_{i}
\label{eq:sbp}
\eeq
which holds without boundary terms in the periodic horizontal directions,
and in the vertical because the boundary contributions carry the factor
$S^u = S^v = 0$ at the lid and bottom faces, the total work of the discrete
force vanishes exactly,
\beq
\sum_{i,j,k} \lk u F^u + v F^v + w F^w \rk = 0
\eeq
for arbitrary $\v u_L$ and arbitrary profiles $u_S(z)$, $v_S(z)$; this has
been verified numerically to machine precision for random fields. Note that
the work vanishes only in the sum.

The vanishing work is a property of the spatial operator and does not by
itself guarantee numerical stability. Since $\v F$ is perpendicular to
$\v u_L$, it turns the velocity vector without changing its magnitude,
i.e.\ it acts locally as a rotation with angular frequency $|\pz{\v u_S}|$,
analogous to inertial oscillations with frequency $f$. An explicit
third-order Adams--Bashforth time stepping scheme is stable for such neutral
oscillations only if the angle covered per time step remains below
$|\pz{\v u_S}| \Delta t \lesssim 0.72$. Since the Stokes shear is largest in
the uppermost grid cell and grows with the developing sea, this number
should be monitored during integrations together with the advective Courant
numbers.

\subsubsection*{Discrete vortex force in the hydrostatic case (pe\_simple)}

In the hydrostatic model pe\_simple the Stokes velocity $\v u_S = (u_S,v_S,0)$
varies in all three dimensions, and the force \eq{eq:vf_cont0} is truncated to
\beq
\v F = \Vec{ - v \, \zeta_S \\ \; u \, \zeta_S }
~~,~~ \zeta_S = \px{v_S} - \py{u_S}
~~,~~ \v u_h^L \cdot \v F = 0
\label{eq:vf_pe}
\eeq
with $\v u_h^L = (u,v)$ in this section for brevity.
The Stokes-shear terms $w \pz{u_S}$ and $w \pz{v_S}$, part of the first and
second component of the total force $-\v u_L \times \vn \times \v u_S$,
respectively, are deliberately omitted. This is because the scalar product
of these terms with $\v u_h^L$ is only cancelled by $w$ times the vertical
force component $-u \pz{u_S} - v \pz{v_S}$, which the hydrostatic model does
not carry. Retained alone, the Stokes-shear terms would thus perform the
uncancelled work $w \lk u \pz{u_S} + v \pz{v_S} \rk$ on $\v u_L$.
The omitted terms are the so-called Stokes shear force of Suzuki and
Fox-Kemper (2016). Their neglect is also in agreement with quasi-geostrophic
scaling where the Rossby number $Ro$ becomes small: within the Stokes layer
of depth $h_S$ the vertical velocity of balanced flow is limited by the
surface kinematic condition and its $O(Ro)$ divergence to
$|w| \lesssim Ro \, U h_S/L$, such that
$|w \pz{u_S}| \sim Ro \, U U_S / L$ independent of $h_S$ and smaller by
$O(Ro)$ than the retained vortex force.

On the C-grid, $u_S$ and $v_S$ are given at the cell centers and averaged
to the $u$- and $v$-points, and the Stokes vorticity is formed at the cell
corners $(i+\frac{1}{2},j+\frac{1}{2})$ by centered differences,
\beq
\zeta_S = \frac{ \Bar{v_S}^{y}_{i+1} - \Bar{v_S}^{y}_{i} }{\Delta x}
        - \frac{ \Bar{u_S}^{x}_{j+1} - \Bar{u_S}^{x}_{j} }{\Delta y}
\eeq
The products are formed at the corners first,
and then averaged back to the velocity points,
\beq
F^u_{i+\frac{1}{2},j} = - \Bar{\zeta_S \, \Bar{v}^{x} }^{y} ~~,~~
F^v_{i,j+\frac{1}{2}} = \; \Bar{ \zeta_S \, \Bar{u}^{y}}^{x}
\label{eq:vf_pe_disc}
\eeq
with $\v F = (F^u,F^v)$.
Applying the summation-by-parts identity \eq{eq:sbp} once in $y$ to the
first and once in $x$ to the second term of the work moves both averages
onto the velocities, and the two contributions meet at the corners,
\beq
\sum u F^u + \sum v F^v
= \sum_{\mathrm{corners}} \zeta_S
 \lk - \Bar{v}^{x} \, \Bar{u}^{y} + \Bar{u}^{y} \, \Bar{v}^{x} \rk = 0
\eeq
i.e.\ the cancellation is not only global as in nh\_simple but holds
corner by corner. The boundary terms of \eq{eq:sbp} vanish because the
velocities vanish on the channel walls and on land, so the identity holds
in the presence of topography; this has been verified to machine precision
for random fields on a masked channel.

\subsubsection*{Discrete aspects of the momentum coupling}

The discretisation of the remaining coupling terms in the ocean model
(C-grid) is chosen such that the energetic properties of the continuous
system are preserved where possible, and monitored where not.

\emph{Stokes--Coriolis force.} Continuously, the combined rotational forces are
$(f+\zeta)\hat z\times\v u_L$, which does no work against $\v u_L$. For the
Coriolis part this property carries over exactly to the discrete level if the
Stokes velocity is first interpolated from T to the U/V points and then passed
through the \emph{same} discrete Coriolis operator $C$ as the Eulerian
velocity (the energy-conserving form of the main code, with $f$ evaluated at
T points inside the adjoint-paired averaging). By linearity the total discrete
Coriolis force is then $C(\v u_L)$, and the antisymmetry of the operator gives
$\sum \v u_L\cdot C(\v u_L)=0$ to machine precision for arbitrary $f(y)$. Note
that adding instead a separate Stokes--Coriolis term with its own (single)
T-to-U/V averaging leaves an $O(\Delta x^2)$ residual in the $\v u_L$ energy
budget, since $\v u_S$ then enters the budget once as a T-point field and once
through a double interpolation; the difference is a commutator between the
$1$-$2$-$1$ smoothing and $f$.

\emph{Vortex force.} With the discretisations of the previous section the
$\v u_L$-work of the applied vortex force vanishes exactly (globally in
nh\_simple, corner by corner in pe\_simple). The coupled driver nevertheless
computes $\sum \v u_L\cdot \v F$ (volume integrated, with the vortex force
exactly as applied) every coupling step and stores it in the coupling
output, so that it can be verified to remain at round-off level.

\emph{Stokes advection of tracers.} The advection of buoyancy by $\v u_S$ is
added in flux form: $\v u_S$ is averaged from T to the U/V faces, the vertical
Stokes velocity $w_S$ is diagnosed from the discrete Stokes continuity
equation integrated upward from the bottom ($w_S(-D)=0$), and the fluxes are
differenced like the model's own tracer fluxes, with zero buoyancy flux
through the surface face. This form is exactly conservative (the fluxes
telescope, so $\int b\,dV$ is preserved to round-off) and
constancy-preserving in the interior (a constant $b$ yields zero tendency,
since the diagnosed $w_S$ closes the three-dimensional Stokes divergence
discretely); both properties were verified to machine precision. In the top
cell the surface no-flux condition places the convergence of the Stokes
transport (the Stokes pumping $w_S$ at the surface face) into the cell budget,
which is the physically correct location: it is compensated by the Eulerian
anti-Stokes flow driven by the Stokes--Coriolis force. Centered fluxes are
used; the Stokes Courant number is $O(10^{-3})$, far below the regime where
flux limiting matters.

\emph{Stress and vertical placement.} The wave stress is kinematic
(the spectrum is the variance density, so no density factor appears; the
air--water density ratio enters through the growth rate of the input source
term) and is applied either to the surface face of the top wet cell
($\tau/\Delta z$) or depth-distributed with the same cell-integrated vertical
kernel that underlies the Stokes velocity and the Doppler projection, which
makes the discrete energy-exchange identity (the $C$-closure of the previous
section) hold to machine precision. Note that with a top-cell thickness large
compared to the Stokes decay scale $1/(2k)$, the stored top-cell Stokes
velocity is a cell average and correspondingly smaller than the surface value;
resolving the near-surface Stokes shear requires $\Delta z$ of a few metres.

\emph{High-frequency tail fluxes to the ocean (tail\_flux switch).} Above
the prognostic cutoff $\nu_c$ the spectrum is diagnostic ($E \propto \nu^{-5}$,
imposed by the tail parameterisation), and its momentum and energy fluxes to
the ocean must be added analytically to the resolved-range dissipation
integrals. Two formulations are implemented, selected by the attribute
\verb|tail_flux|:

(i) \verb|'wamdi'| (default): the WAMDI dissipation rate
continued over the tail,
$\int_{\nu_c}^\infty \gamma_{out}\,E_{hi}\, \sigma\, d\nu$ with $\gamma_{out} \propto
\bar\sigma (\sigma/\bar\sigma)^2 (\bar\alpha/\alpha_{PM})^2$. The integrand
decays like $\nu^{-2}$ and converges; note that the corresponding integral with
the dissipation of Janssen (1994) diverges strongly because of its
$(k/\bar k)^4$ term (its $(k/\bar k)^2$ term behaves like the WAMDI form),
which is why the WAMDI form is used for this term. However,
$\gamma_{out}$ is the whitecapping rate calibrated on the resolved spectrum and
underestimates the dissipation that actually maintains the $\nu^{-5}$
equilibrium range --- the very reason the tail must be imposed diagnostically
--- so this form systematically underestimates the tail fluxes (in tests by
a substantial factor in the momentum flux).

(ii) \verb|'input'|: in the equilibrium range $S_{in} + S_{nl} + S_{out} = 0$
holds, so the flux to the ocean, $-S_{out}$, equals $S_{in} + S_{nl}$. The
input part is the analytic integral of the WAMDI growth rate over the tail:
with $a(\theta) = 28 u_\star 2\pi \cos(\theta-\theta_w)/g$ (the input
parameter evaluated at $\nu_c$), the input being active only above
$\nu_1 = 1/a$,
\beq
I_\tau(\theta) = \frac{\epsilon}{4} (2\pi)^2 \nu_c^3\, E_c(\theta)
\left\{ \begin{array}{ll} a \nu_c - \frac{1}{2}, & a \nu_c \ge 1\\[2pt]
\frac{1}{2} (a \nu_c)^2, & 0 < a \nu_c < 1 \end{array} \right.
\qquad
I_\Phi(\theta) = \frac{\epsilon}{4} (2\pi) \nu_c^2 g\, E_c(\theta)
\left\{ \begin{array}{ll} \frac{a \nu_c}{2} - \frac{1}{3}, & a \nu_c \ge 1\\[2pt]
\frac{1}{6} (a \nu_c)^3, & 0 < a \nu_c < 1 \end{array} \right.
\eeq
($\epsilon = \rho_a/\rho_w$, $E_c(\theta) = E(\nu_c,\theta)$; continuous at
$a \nu_c = 1$; zero upwind; the integrands again decay like $\nu^{-2}$). The same booking applies to the
\emph{in-grid} diagnostic bins (between the per-column prognostic cut-off
$m_0$ and the last grid frequency, where the $\nu^{-5}$ shape is imposed):
there the flux is evaluated discretely as $\gamma_{in} E$ with the actual
imposed spectrum, while the prognostic bins keep the true $\gamma_{out} E$
dissipation --- booking the in-grid diagnostic range at $\gamma_{out}$ would be
inconsistent with the equilibrium assertion that the imposed tail itself
encodes (and underestimates the flux substantially: the in-grid diagnostic
range can carry several times the beyond-grid tail). The nonlinear part uses the
conservation of the exact transfer over the full spectrum,
$\int_{tail} S_{nl}\, d\Omega = -\int_{resolved} S_{nl}\, d\Omega$: the
\emph{prognostic}-range integrals of the stored $S_{nl}$ (momentum weight
$\sigma$, energy weight $g$) are credited to the ocean fluxes (the ``cascade
credit''; the in-grid diagnostic $S_{nl}$ is part of the equilibrium booking
there, and the algebra of the two ranges collapses exactly to the
prognostic-masked integral).
On the truncated grid the discrete DIA is not exactly conservative at the
cutoff; the residual is deliberately included in the credit --- whatever the
discrete operator removes from the resolved range is deposited in the tail.
For the depth-distributed stress (method \verb|calc_stress_work|, as opposed
to the surface stress of \verb|calc_stress|) the tail is
deposited with its exact profile: with $b = 8\pi^2/g$, $t=-z$,
$\beta_L = b\,t \nu_L^2$, the cumulative-from-surface tail stress is
$V(t) = C_0 [\, a/\nu_L - 1/(2\nu_L^2) - (a J_2 - J_3)\,]$ with
$J_2 = e^{-\beta_L}/\nu_L - \sqrt{\pi b t}\,\mathrm{erfc}\sqrt{\beta_L}$ and
$J_3 = \frac{1}{2}( e^{-\beta_L}/\nu_L^2 - b t\, E_1(\beta_L))$, so the cell
weights $(V(t_{bot})-V(t_{top}))/\Delta z$ sum to the surface tail exactly;
the cascade credit is deposited in the uppermost cell (its carrier waves
have deposition scales $g/2\sigma^2$ of centimetres).

\emph{Dissipation energy flux to the ocean.} The term
$\int \sigma S_{out}\, d\v k$ of the wave energy budget \eq{eq_W} is
the energy the breaking waves hand to near-surface turbulence. It is diagnosed
locally as
\beq
\Phi_{oc}(\v x) = -\,g \int \dot E_{out}\, d\nu\, d\theta \;\ge\, 0
\eeq
(kinematic; $\rho_w \Phi_{oc}$ in W\,m$^{-2}$; WAM's PHIOC), computed from the
same stored dissipation rate $\gamma_{out}$ and with the same spectral
measures as the wave stress, and written to the coupling output both as the
local field and as its domain integral in the energy-diagnostic record. The
dissipation of the diagnostic $\nu^{-5}$ tail is included according to the
\verb|tail_flux| switch described above (under \verb|'input'| together with
the nonlinear cascade credit). The
field is the surface boundary condition for the flux $F$ of the turbulent
kinetic energy (TKE) closure following Gaspar et al.\ (1990) in the ocean
component ($\rho_w \Phi_{oc}$ as the surface TKE flux, in place of the
$u_\star^3$-based parameterisation), which routes this term into the ocean's
turbulent kinetic energy budget and thereby closes the coupled energy cycle.
The TKE closure receives, besides the usual shear production by the Eulerian
mean flow, a shear production term related to the mean Stokes velocity,
$c_u K (\p_z \bar{\v u}) \cdot \p_z \bar{\v u}_S$, resulting from the
down-gradient closure of the eddy stresses; the corresponding energy is
exchanged with the mean Lagrangian kinetic energy.

\emph{Vertical friction and mixing.} For consistency with the TKE closure,
the vertical friction operator acts on the Eulerian mean velocity,
$\p_z A_v \p_z (\v u_L - \v u_S) = \p_z A_v \p_z \v u$: the work of this
viscous term against $\v u_L$ then exactly matches the shear production
terms of the TKE closure, including the Stokes-shear production. Mixing the
Lagrangian instead of the Eulerian momentum would break this energy
consistency.

\emph{Coupling frequency and the exchange term.} The two models exchange
fields once per coupling interval $\Delta t_c$, chosen as a common multiple of
the two time steps: the component with the smaller time step sub-cycles
$\Delta t_c/\Delta t$ steps between exchanges. When the wave step is the
smaller one, $\Delta t_c = \Delta t_{\rm ocean}$ and the wave model
sub-cycles; when $\Delta t_{\rm ocean} < \Delta t_{\rm wave}$ (e.g.\ a
non-hydrostatic or very high resolution ocean) the roles reverse and the ocean
sub-cycles, with $\Delta t_c = \Delta t_{\rm wave}$. In either case all
coupling fields are piecewise constant over the coupling interval, so the
exchange work $C=\int \v u_L\cdot\p_t\v u_S\,dV$ is evaluated with
$\p_t\v u_S$ as the difference quotient across the coupling update, i.e.\ over
$\Delta t_c$ rather than over the time step of either component. The
domain-integrated energy diagnostics --- wind input to the waves
$g\int\gamma_{\rm in}E\,d\Omega\,dA$, the stress work on the ocean (computed
with the identical stencils as the applied forcing), the exchange term $C$,
and the vortex-force work --- are written to the coupling output at every
coupling step, so the discrete energy pathway of the coupled system is
monitored throughout the integration rather than assumed. The choice of
$\Delta t_c$ is constrained by the physics rather than by stability: the
coupling fields should be well sampled, i.e.\ $\Delta t_c$ small compared to
the evolution time scales of the exchanged quantities (translation of the wind
systems across a grid cell for the stress and Stokes fields, and the evolution
of the resolved ocean currents for the refraction terms).

  
\subsection*{Propagation, boundary conditions, land mask}

Action is propagated since it avoids to calculate the energy exchange with the mean flow
explicitely. 
Advection uses the second-order scheme with Superbee limiter (\verb|tracer_flux_superbee|)
taken from MITgcm. 
A first-order upwind option is also available (\verb|tracer_flux_upwind|,
selected by the attribute ''f\_limiter''$=$''upwind''); it shares the identical
finite-volume (face-flux) structure --- the same land masking, halo exchange,
closed frequency faces and reflection --- so conservation and the energy fixer
of \eq{eq:efix} still apply, but the flux limiter is set to zero. Upwind is more
diffusive (it smears sharp swell fronts) but robust to grid-scale noise.
The refraction drift velocities $\dot \nu$ and $\dot \theta$ can optionally be
clipped to a CFL bound (attribute \verb|limit_fdot|, default off, with the
bound set by \verb|refraction_cfl|, default $0.5$): strong current shear can
otherwise drive the spectral advection beyond its stability limit locally.
The clipping breaks the exact ray equations where it engages and is meant as
a safeguard, not a substitute for a proper time step.



Lateral boundary conditions in $x$ and $y$ are periodic (attribute ''periodic\_in\_x''/''periodic\_in\_y''),
closed, or reflective. Boundary conditions in frequency are always closed, in direction
always periodic.
Lateral 
reflection boundary conditions at coast lines are implemented  the same way as in the 
IDEMIX2 in pyOM2.
These are activated by setting the attribute ''enable\_reflection'' to true value'. 
If undefined or false, the default of closed boundary are used where energy may pile up near the 
boundary.
Land mask is set in the template method   ''set\_topography'' in \verb|./pyOM3/wim/main.py|.


\subsection*{Time stepping scheme}

Each time step is split into an explicit propagation substep for the wave action
$A$ followed by an implicit source substep for the wave energy $E = \sigma A$.
The propagation tendency $\dot A|_{\rm prop}$ (the sum of the flux divergences in
$x$, $y$, $\nu$ and $\theta$) is advanced with a multi-step Adams--Bashforth
scheme,
\beq \label{eq:ab3}
A^{n+1} = A^{n} + \Delta t\left( c_A\, \dot A^{\,n}|_{\rm prop}
                              + c_B\, \dot A^{\,n-1}|_{\rm prop}
                              + c_C\, \dot A^{\,n-2}|_{\rm prop}\right).
\eeq
The scheme used is the stabilised second- or third-order Adams--Bashforth  pair
\beq \label{eq:ab2eps}
(c_A,c_B,c_C) &=& \left(\frac{3}{2}+\epsilon,\; -\frac{1}{2}-\epsilon,\; 0\right),
\qquad \epsilon \approx 0.05, ~~\mbox{for second order}
\\
(c_A,c_B,c_C) &=& \left(\frac{23}{12},\; -\frac{16}{12},\; \frac{5}{12}\right)
~~\mbox{for third order}
\eeq
both started with a forward Euler step $(1,0,0)$. The time loop uses the
AB2$+\epsilon$ pair; the third-order coefficients are given for reference but
are deliberately not used (see below). The recent tendencies are held in a
rotating three-level buffer (the pointers ''tau'', ''taum1'', ''taum2''). Only
the propagation is integrated this way; the source terms are treated separately
below. After the propagation update the action is subject to the lateral
boundary conditions and, on land, is masked to zero.

The choice of the time scheme for the propagation is more delicate than usual
because the advection uses a flux \emph{limiter}. TVD limiters (Superbee,
upwind) guarantee a monotone, non-oscillatory update --- in particular
$A \ge 0$ --- only when the limited flux is applied in a \emph{single forward
Euler step} (or a strong-stability-preserving Runge--Kutta combination of such
steps). Multi-step Adams--Bashforth schemes combine tendencies from several time
levels with \emph{negative} weights, which destroys this guarantee at any CFL
number: each stored tendency is limited with respect to its own state, but the
negatively-weighted combination is free to over- and undershoot. For the
third-order scheme $(23/12,\,-16/12,\,5/12)$ the negative weight is large and
the violation is severe: a smooth, sign-definite action field develops
undershoots ($A < 0$) in an alternating grid-scale ($2\Delta x$) checkerboard
pattern at sharp fronts (storm edges, swell fronts). Unlike in the ocean
model --- where the same combination is harmless because tracers are not
sign-definite, explicit (biharmonic and vertical) diffusion damps grid-scale
noise every step, and no term amplifies small wiggles --- the wave model has no
explicit diffusion, its action must remain non-negative, and the source terms
amplify any $\pm$ pattern strongly non-linearly (through
$\alpha = \bar E\,\bar\sigma^4$ and the cubic DIA transfer), so the checkerboard
grows over $\mathcal{O}(1)$ day until floating-point overflow. The AB2$+\epsilon$
coefficients \eq{eq:ab2eps} reduce the negative weight to $-(1/2+\epsilon)$,
which in practice (CFL${}\lesssim 0.25$) is small enough for the limiter's
built-in slack to absorb the violation: the action remains non-negative to
machine precision in long integrations where AB3 fails within a day. This is an
empirical rather than a rigorous property; forward Euler $(1,0,0)$ or an SSP
Runge--Kutta scheme restores the formal TVD guarantee at the cost of temporal
accuracy or an additional flux evaluation per step.

Two safeguards back this up. First, a positivity floor (attribute
\verb|floor_action|, default on) clips the action at zero after each
propagation step, as is standard in WAM. The clipped mass --- the action the
floor removes, which equals the spurious action it injects and hence the
energy-consistency cost of the clip --- is measured \emph{before} clipping and
a warning is issued whenever it exceeds round-off scale relative to the total
action, so a silently active floor cannot mask a genuine positivity violation;
switching the floor off restores the raw scheme. Second, $\min A$ is monitored
automatically every $O(100)$ steps against a relative round-off threshold and a
warning names the remedy (larger $\epsilon$ or forward Euler): any
limiter/time-scheme violation shows up there long before a blow-up.


Writing the source terms as growth rates acting on the energy, the source substep
integrates $\partial_t E = S$ with $S = \dot E_{nl} + (\gamma_{in} +
\gamma_{out})E$, where $\dot E_{nl}$ is the non-linear (DIA) transfer. A fully
explicit update would require a very small time step because the growth rate
increases with frequency (this is also the reason for the high-frequency
cut-off). Instead a linearised-implicit (semi-implicit) step is used: linearising
$S$ about the current state, $S(E^{n+1}) \approx S(E^n) + \partial_E S\,
(E^{n+1}-E^n)$, and solving for $E^{n+1}$ gives
\beq \label{eq:implicit_source}
E^{n+1} = E^{n} + \frac{S\, \Delta t}{1 - \Delta t\, \partial_E S},
\qquad
\partial_E S = \gamma_{in} + \gamma_{out} + \partial_E \dot E_{nl},
\eeq
where the diagonal of the non-linear functional derivative
$\partial_E \dot E_{nl}$ is returned by the DIA backend alongside the transfer
itself. This is unconditionally stable for a purely dissipative feedback
($\partial_E S < 0$), for which the denominator exceeds unity and the update is
damped. The growth can additionally be bounded as in WAM by a growth limiter
(attribute \verb|limit_growth|, default off) that caps $S$ (not $\partial_E S$)
so that the energy cannot grow by more than a prescribed fraction per step,
\beq \label{eq:growthlimit}
S \leftarrow \min\!\left(S,\; 10^{-7}\, g\,
\max(u_\star,\,u_\star^{\min})\, \bar\sigma\,/\,(2\pi \nu^4)\right),
\qquad u_\star^{\min} = 0.01\ {\rm m\,s^{-1}}.
\eeq
The friction-velocity floor $u_\star^{\min}$ in the cap is essential and easy
to overlook: without it, the cap vanishes wherever the wind vanishes and the
limiter degenerates to $S \leftarrow \min(S,0)$ in calm cells. This clips only
the \emph{positive} lobes of the nonlinear transfer, destroying the
conservative $\pm$ structure of $S_{nl}$ exactly there, and acts as a small
but systematic negative-action source. In a calm region that accumulates swell
without a sink (e.g.\ a walled basin with reflective coasts outside the wind
band), the resulting violation grows linearly in time --- independent of the
time step and of the time-stepping scheme --- until, after $O(100)$ days, the
nonlinear source terms of the accumulated spectrum activate and the
integration fails abruptly. The floor keeps the cap strictly positive
everywhere while leaving the physics unchanged: $u_\star^{\min}$ is far below
the input threshold of the slowest resolved wave ($c \approx 28\,u_\star$
requires $u_\star \gtrsim 0.06$ m\,s$^{-1}$ for the highest resolved
frequency), so the wind input remains identically zero in floored cells; only
the growth cap becomes a finite positive bound instead of zero.
Because the limiter bounds only the numerator, the denominator
$1 - \Delta t\, \partial_E S$ can still vanish or change sign where
$\partial_E S$ is large and positive --- most notably at high frequency during
spin-up from a small seed, where the DIA derivative is steep. When this happens
the update overshoots and the highest bins become unstable. Two safeguards are
used: the high-frequency cut-off removes the steepest bins from the prognostic
integration (replacing them by the diagnostic $\nu^{-5}$ tail), and, optionally,
only the dissipative part of the feedback is kept in the implicit denominator,
\beq \label{eq:denomclamp}
\partial_E S \;\rightarrow\; \min\!\left(\partial_E S,\, 0\right),
\eeq
guaranteeing $1 - \Delta t\, \partial_E S \ge 1$. This is activated 
by setting the attribute  \verb|implicit_dissipation_only|  to
\verb|True| and equals  the standard WAM/WW3
construction --- the stabilising dissipative feedback is treated implicitly while
the growth is treated explicitly and bounded by the limiter. It leaves the
equilibrium ($S = 0$) unchanged and therefore does not affect the energy balance;
it only alters the transient integration of the growth phase.

After the source substep the energy is masked to zero on land, the
high-frequency tail is re-imposed through \verb|imphftail|, and the action is
recovered as $A = E/\sigma$ for the next propagation step.


The time step of the propagation is limited by the advective CFL number based
on the fastest group velocity, i.e.\ the lowest resolved frequency (Superbee is
stable to CFL${}\lesssim 0.5$, upwind to CFL${}<1$; in practice CFL${}\approx
0.2$--$0.25$ is used, which is also the regime in which the AB2$+\epsilon$
scheme preserves the non-negativity of the action, see the time-stepping
section). Note that the grid-scale checkerboard described there is \emph{not} a
CFL violation --- it appears at any Courant number when a multi-step scheme
with a large negative weight (AB3) is combined with the flux limiter --- so it
cannot be cured by reducing the time step; the cure is the time scheme.
Substantially larger time steps than the advective CFL allows would require an
implicit or semi-Lagrangian treatment of propagation, at the cost --- for plain
semi-Lagrangian --- of the exact (face-flux) energy conservation.


\subsection*{Model parameters and options}

All parameters and options are attributes of the model instance. They are set
either through the parameter dictionary passed at construction or by assigning
the attribute on the instance before the model is stepped; options marked with
a default may simply be left unset. The tables below list every attribute the
wave model reads, with its default where one exists.

\subsubsection*{Required parameters (no default)}

\begin{longtable}{@{}llp{8.6cm}@{}}
\hline
Attribute & Unit & Meaning \\
\hline
\endhead
\verb|nx|, \verb|ny| & --- & global horizontal grid size (cells, without halo
  points) \\
\verb|dx|, \verb|dy| & m & horizontal grid spacing \\
\verb|nf| & --- & number of frequency bins (local on every PE) \\
\verb|nd| & --- & number of direction bins (local on every PE) \\
\verb|f0| & 1/s & lowest frequency of the geometric frequency grid (WAM value
  $0.035$) \\
\verb|df0| & --- & relative frequency increment $\Delta \nu/\nu$ of the
  geometric grid (WAM value $0.1$) \\
\verb|dt| & s & time step \\
\verb|periodic_in_x| & --- & lateral boundary in $x$: periodic if true, closed
  if false (the outermost ring of cells is then set to land) \\
\verb|periodic_in_y| & --- & as above for $y$ \\
\hline
\end{longtable}

\noindent
The constructor additionally takes the keyword arguments \verb|n_pes_x|,
\verb|n_pes_y| (processor grid, default $1\times 1$; their product must equal
the communicator size), \verb|halo_points| (default 1) and \verb|comm_in|
(default: the MPI world communicator).

\subsubsection*{Forcing and coupling fields}

\begin{longtable}{@{}llp{8.6cm}@{}}
\hline
Attribute & Default & Meaning \\
\hline
\endhead
\verb|ustar| & required & friction velocity magnitude $u_\star$; scalar or
  horizontal field (scalars are broadcast) \\
\verb|udir| & required & wind direction $\theta^*$ (radians); scalar or field \\
\verb|ocean_U|, \verb|ocean_V| & unset & Doppler velocity $\v U(\nu)$ provided
  by the coupler. If unset, all current effects (advection by currents,
  refraction, the energy fixer) are dormant \\
\verb|ocean_Udf|, \verb|ocean_Vdf| & $0$ & frequency derivative
  $\p_\nu \v U$ of the Doppler velocity (radial group-velocity correction
  $R$ and frequency refraction) \\
\hline
\end{longtable}

\subsubsection*{Options with defaults}

\begin{longtable}{@{}llp{8.2cm}@{}}
\hline
Attribute & Default & Meaning \\
\hline
\endhead
\verb|snonlin_backend| & \verb|'fortran'| & DIA implementation:
  \verb|'fortran'| (original WAM code; requires double precision),
  \verb|'jax'| or \verb|'numpy'|; all give identical results \\
\verb|grav| & $9.81$ & gravitational acceleration \\
\verb|rho_o| & $1024$ & reference water density. Currently unused inside the
  wave model: the air--water density ratio of the source terms is fixed at
  $1.225/1000$ \\
\verb|z0| & $0.0002$ & roughness length; read only by the optional WAM input
  parameterisation, inactive with the default WAMDI input \\
\verb|mask_energy| & off & prognostic cut-off $\nu_c$ per column with the
  imposed $\nu^{-5}$ diagnostic tail (\verb|imphftail|); prerequisite for the
  equilibrium booking of \verb|tail_flux| $=$ \verb|'input'| \\
\verb|limit_growth| & off & WAM growth limiter, \eq{eq:growthlimit} \\
\verb|implicit_dissipation_only| & off & keep only the dissipative part of the
  source feedback in the implicit denominator, \eq{eq:denomclamp} \\
\verb|floor_action| & on & positivity floor $A \leftarrow \max(A,0)$ after the
  propagation step, with the clipped mass monitored and warned about \\
\verb|f_limiter| & \verb|'superbee'| & advection flux limiter:
  \verb|'superbee'| (second order) or \verb|'upwind'| (first order,
  diffusive, robust) \\
\verb|energy_consistent| & off & conservative energy fixer of the propagation
  operator, \eq{eq:efix} \\
\verb|enable_reflection| & off & specular reflection of the advective fluxes
  at coastlines \\
\verb|limit_fdot| & off & clip the refraction velocities $\dot\nu, \dot\theta$
  to a CFL bound \\
\verb|refraction_cfl| & $0.5$ & the CFL bound used by \verb|limit_fdot| \\
\verb|closed_bottom| & off & extend the bottom cell of the coupling kernel to
  $z \to -\infty$ (exact barotropic projection, no momentum lost below the
  column floor) \\
\verb|tail_flux| & \verb|'wamdi'| & booking of the diagnostic-tail fluxes to
  the ocean: \verb|'wamdi'| (dissipation-rate continuation) or
  \verb|'input'| (equilibrium-range $\gamma_{in}$ integral plus nonlinear
  cascade credit) \\
\hline
\end{longtable}

\subsubsection*{Snapshot diagnostics}

The netCDF snapshot class reads the attribute \verb|snapint| (snapshot
interval in seconds, required) and takes the setup arguments
\verb|snap_file_name| (default \verb|'snap.cdf'|) and
\verb|show_each_time_step| (default on). The land mask is defined by
overloading the template method \verb|set_topography|; the default is an
all-ocean domain, with closed outer boundaries generated automatically from
\verb|periodic_in_x|/\verb|periodic_in_y|.


\end{document}
